\section{Étape 5 -- Vérification de la configuration}
\label{secEtapeVerification}

Avant tout téléchargement dans l'automate, la configuration établie doit
être vérifiée.

\subsection{Sauvegarde et compilation}

Sauvegarder l'application, puis la compiler.

\subsection{Contrôle de la cartographie mémoire}

À l'aide de l'éditeur de données (variables dérivées), vérifier que les
zones mémoire allouées à l'équipement ajouté correspondent à l'attendu,
c'est-à-dire :

\begin{itemize}
    \item pour les données produites (\texttt{T->O}) : une variable du type
    utilisateur créé pour l'équipement (par exemple \texttt{guichet\_IN} de
    type \texttt{T\_guichet\_IN}), positionnée à l'adresse mémoire attendue
    d'après la cartographie établie pour le coupleur ;
    \item pour les données consommées (\texttt{O->T}) : de même, une
    variable du type utilisateur correspondant (par exemple
    \texttt{guichet\_OUT} de type \texttt{T\_guichet\_OUT}), à l'adresse
    mémoire attendue.
\end{itemize}

\begin{UPSTIwarning}{Validation avant téléchargement}
La configuration réalisée doit être vérifiée (par exemple par l'enseignant,
en contexte pédagogique, ou par un pair en contexte professionnel) avant
tout téléchargement dans l'automate. Une erreur de cartographie mémoire ou
d'adressage IP peut en effet perturber le fonctionnement d'autres
équipements présents sur le même réseau.
\end{UPSTIwarning}

\subsection{Configuration de la tâche}

Une fois la configuration réseau validée, la tâche périodique de
l'automate (dossier \texttt{MAST}) doit être paramétrée pour tenir compte
du procédé piloté : période de scrutation et surveillance du temps de
cycle (chien de garde), dont les valeurs sont déterminées à partir de
l'analyse de la dynamique du procédé (cf. document
\texttt{02b\_donnees\_echangees} pour la description des données
échangées).
